% Source: https://artofproblemsolving.com/community/c3h1691857_prove_that_division_is_repeated_subtraction
% Stable identity: aops-topic-1691857-q01
After $k$ subtractions the result is $a-kb=(q-k)b$. This is positive for $k<q$ and equal to $0$ for $k=q$, so the result first reaches $0$ after exactly $q$ subtractions. Hence $a\div b$ is the \emph{number of times} $b$ can be subtracted from $a$ until nothing is left.

In the example, $4-2=2$ and $2-2=0$: two subtractions of $2$ empty the $4$, so $4\div2=\boxed{2}$. The $0$ at the end is what is left over (the remainder), not the quotient; the quotient is the count of subtractions. (In general, for $a=qb+r$ with $0\le r<b$, subtracting $b$ until the result is less than $b$ takes $q$ steps and leaves $r$.)
